\documentclass{book}
%\documentclass{book}
%If you're going to print this, perhaps get ride of oneside so that \parts{} will have their own page!
\title{
	\begin{huge}
		\bf{Hirzebruch–Riemann–Roch theorem}
	\end{huge}
	}
\author{Nathanael Chwojko-Srawley}

%\date{} % If you want to specify a date

%  ╔══════════════════════════════════════════════════════════╗
%  ║                         packages                         ║
%  ╚══════════════════════════════════════════════════════════╝

\usepackage[a4paper, total={6in, 8in}]{geometry}
\usepackage[answerdelayed]{exercise}
\usepackage[font=itshape]{quoting}
\usepackage[most]{tcolorbox}         % Makes nice boxes for thms, cor, prop, or anything else.
\usepackage{amsmath, amssymb, amsthm} % Used to write better mathematical expressions (ex. \mathbb{R})
\usepackage{array}                   % \begin{table}{>{\em}c c} -> 1st column italic (bfseries for bold)
\usepackage{bbm}                     % for mathbbm{1}
\usepackage{blindtext}
\usepackage{commutative-diagrams}
\usepackage{dsfont}
\usepackage{dynkin-diagrams}         % for Dynkin diagrams
\usepackage{enumitem}
\usepackage{etoolbox}
\usepackage{euscript}
\usepackage{extarrows}               % Used to compile any UTF-8 Character (and supports Unicode)
\usepackage{fancyhdr}
\usepackage{float}                   % package with ``H'' for figures and tables, ex. Images, tables, etc
\usepackage{graphicx}                % Let's you use images in LaTeX; ex the \includegraphics command.
\usepackage{hhline}
\usepackage{hyperref}
\usepackage{cleveref}
\usepackage{inputenc}                % Used to compile any UTF-8 Character (and supports Unicode)
\usepackage{listings}                % Create nice colour code blocks (customized in preamble)
\usepackage{longtable}               % create tables that extend past one page
\usepackage{makeidx}
\usepackage{mathrsfs}
\usepackage{mathtools}
% \usepackage{microtype}               % makes nice text. Fixes some under-/over- filled box problems.
\usepackage{multicol}                % For multiple columns
\usepackage{pgfplots}
\usepackage{scalerel}
\usepackage{setspace}                % More flexibility on spacing in document.
\usepackage{stackengine}
\usepackage{stackrel}                % for left-right arrows, staking symbol above and bellow
\usepackage{thmtools}                % Customize theorem environment (in preamble).
\usepackage{tikz-cd}                 % for commutative diagrams https://tex.stackexchange.com/questions/419221/how-to-do-the-pushout-with-universal-property
\usepackage{tikz}                    % for commutative diagrams https://tex.stackexchange.com/questions/419221/how-to-do-the-pushout-with-universal-property
\usepackage{titlesec}
\usepackage{wasysym}                 % emoji's!
\usepackage{wrapfig}                 % To wrap text around figure
\usepackage{xcolor}                  % Colour text.

% \usepackage{comment}

% \usepackage{CJKutf8} % for some chinese

% \usepackage{dutchcal} % for lower case mathcal
% \pdfsuppresswarningpagegroup=1

% ╔═════════════════════════════════════════════════════╗
% ║           for figures via inkspace                  ║
% ╚═════════════════════════════════════════════════════╝

% to import figures via: https://github.com/gillescastel/inkscape-figures
\usepackage{import}
\usepackage{pdfpages}
\usepackage{transparent}


% This command is the ONLY command imported via the packages snippet, think of it as a tiny package written out
\newcommand{\incfig}[2][1]{%
    \def\svgwidth{#1\columnwidth}
    \import{./figures/}{#2.pdf_tex}
}

% ╔═══════════════════════════════════════════════════════════╗
% ║                 bibliography and citing                   ║
% ╚═══════════════════════════════════════════════════════════╝

\usepackage{csquotes}
% \usepackage[backend=biber,citestyle=alphabetic,bibstyle=authortitle]{biblatex}
\usepackage[backend=biber]{biblatex}

\addbibresource{bibliography.bib}

% ╔═══════════════════════════════════════════════════════════╗
% ║                    colours (colors)                       ║
% ╚═══════════════════════════════════════════════════════════╝

\definecolor{dkgreen}{rgb}{0,0.6,0}
\definecolor{gray}{rgb}{0.5,0.5,0.5}
\definecolor{mauve}{rgb}{0.58,0,0.82}
\definecolor{lightyellow}{rgb}{1,1,0.6}
\definecolor{reddish}{HTML}{A01A3A}

%  ╔══════════════════════════════════════════════════════════╗
%  ║               title/heading/footer format                ║
%  ╚══════════════════════════════════════════════════════════╝

\pagestyle{fancy}
\setlength\headheight{20pt}
\renewcommand{\sectionmark}[1]{\markright{\thesection.\ #1}}
\renewcommand{\chaptermark}[1]{\markboth{#1}{}}
\fancyfoot[C,CO]{\textcolor{gray}{Nathanael Chwojko-Srawley}}
\fancyfoot[R,RO]{\thepage}{}

\titleformat
{\chapter} % command
[display] % shape
{\bfseries\Huge\itshape} % format
{\thechapter} % label
{0.5ex} % sep
{
    \rule{\textwidth}{1pt}
    \vspace{1ex}
    \centering
}
[
\vspace{-0.5ex}
\rule{\textwidth}{0.3pt}
]

\setlength{\parindent}{0pt} % don't like indentation infront of paragraphs
\setlength{\parskip}{1ex}   %so that there is still space between paragarphs


%  ╔══════════════════════════════════════════════════════════╗
%  ║                       Shortcuts                          ║
%  ╚══════════════════════════════════════════════════════════╝

%\ExerciseLevelInToc

% I like original mathcal better, but need lowercase.
\DeclareFontFamily{OT1}{pzc}{}
\DeclareFontShape{OT1}{pzc}{m}{it}{<-> s * [1.10] pzcmi7t}{}
\DeclareMathAlphabet{\mathpzc}{OT1}{pzc}{m}{it}


% mathbb
\newcommand{\A}{{\mathbb{A}}}
\newcommand{\C}{{\mathbb{C}}}
\newcommand{\F}{{\mathbb{F}}}
\newcommand{\N}{{\mathbb{N}}}
\newcommand{\Q}{{\mathbb{Q}}}
\newcommand{\R}{{\mathbb{R}}}
\newcommand{\V}{{\mathbb{V}}}
\newcommand{\Z}{{\mathbb{Z}}}
\newcommand{\BI}{{\mathbb{I}}}
\renewcommand{\H}{{\mathbb{H}}}
\renewcommand{\P}{{\mathbb{P}}}


% mathcal
% \newcommand{\co}{{\mathfrak{o}}} %mathcal{o} looks like wreath product, so I changed it to mathfrak
\newcommand{\co}{{\mathpzc{o}}}
\newcommand{\CA}{{\mathcal{A}}}
\newcommand{\CB}{{\mathcal{B}}}
\newcommand{\CC}{{\mathcal{C}}}
\newcommand{\CD}{{\mathcal{D}}}
\newcommand{\CF}{{\mathcal{F}}}
\newcommand{\CG}{{\mathcal{G}}}
\newcommand{\CH}{{\mathcal{H}}}
\newcommand{\CI}{{\mathcal{I}}}
\newcommand{\CK}{{\mathcal{K}}}
\newcommand{\CL}{{\mathcal{L}}}
\newcommand{\CM}{{\mathcal{M}}}
\newcommand{\CN}{{\mathcal{N}}}
\newcommand{\CO}{{\mathcal{O}}}
\newcommand{\CQ}{{\mathcal{Q}}}
\newcommand{\CR}{{\mathcal{R}}}
\newcommand{\CS}{{\mathcal{S}}}
\newcommand{\CT}{{\mathcal{T}}}
\newcommand{\CV}{{\mathcal{V}}}
\newcommand{\CZ}{{\mathcal{Z}}}
% EXCPETION:
% \newcommand{\CP}{{\mathcal{P}}}
\newcommand{\CP}{\mathbb{C}\mathrm{P}}
\newcommand{\RP}{\mathbb{R}\mathrm{P}}


%Euscript
\newcommand{\EB}{\EuScript{B}}
\newcommand{\EC}{{\EuScript{C}}}
\newcommand{\EF}{{\EuScript{F}}}
\newcommand{\EL}{{\EuScript{L}}}
\newcommand{\EO}{{\EuScript{O}}}


%mathfrak (lowercase)
\newcommand{\fa}{{\mathfrak{a}}}
\newcommand{\fb}{{\mathfrak{b}}}
\newcommand{\fc}{{\mathfrak{c}}}
\newcommand{\fd}{{\mathfrak{d}}}
\newcommand{\fm}{{\mathfrak{m}}}
\newcommand{\fn}{{\mathfrak{n}}}
\newcommand{\fo}{{\mathfrak{o}}}
\newcommand{\fp}{{\mathfrak{p}}}
\newcommand{\fq}{{\mathfrak{q}}}


%mathfrak (upper case)
\newcommand{\FA}{{\mathfrak{A}}}
\newcommand{\FB}{{\mathfrak{B}}}
\newcommand{\FC}{{\mathfrak{C}}}
\newcommand{\FD}{{\mathfrak{D}}}
\newcommand{\FK}{{\mathfrak{K}}}
\newcommand{\FM}{{\mathfrak{M}}}
\newcommand{\FN}{{\mathfrak{N}}}
\newcommand{\FO}{{\mathfrak{O}}}
\newcommand{\FP}{{\mathfrak{P}}}


%mathscr
\newcommand{\ff}{{\mathscr{F}}}
\newcommand{\SA}{\mathscr{A}}
\newcommand{\SC}{\mathscr{C}}
\newcommand{\SF}{\mathscr{F}}
\newcommand{\SG}{\mathscr{G}}
\newcommand{\SH}{\mathscr{H}}
\newcommand{\SI}{\mathscr{I}}
\newcommand{\SK}{\mathscr{K}}
\newcommand{\SN}{\mathscr{N}}
\newcommand{\SQ}{\mathscr{Q}}
\newcommand{\SR}{\mathscr{R}}
\newcommand{\SV}{\mathscr{V}}
\newcommand{\SZ}{\mathscr{Z}}


% operatorname -- 2 letters (lowercase and upper case)
\newcommand{\ab}{{\operatorname{ab}}}
\newcommand{\alg}{{\operatorname{\alg}}}
\newcommand{\ad}{{\operatorname{ad}}}
\newcommand{\ch}{{\operatorname{ch}}}
\newcommand{\ev}{{\operatorname{ev}}}
\newcommand{\id}{{\operatorname{id}}}
\newcommand{\im}{{\operatorname{im}}}
\newcommand{\nm}{{\operatorname{Nm}}}
\newcommand{\op}{{\operatorname{op}}}
\newcommand{\rk}{{\operatorname{rk}}}
\newcommand{\sh}{{\operatorname{sh}}}
\newcommand{\tr}{{\operatorname{tr}}}

\newcommand{\Ab}{{\operatorname{Ab}}}
\newcommand{\Br}{{\operatorname{Br}}}
\newcommand{\Ch}{{\operatorname{Ch}}}
\newcommand{\Cl}{{\operatorname{Cl}}}
\newcommand{\GL}{{\operatorname{GL}}}
\newcommand{\Nm}{{\operatorname{Nm}}}
\newcommand{\SL}{{\operatorname{SL}}}

\renewcommand{\Re}{{\operatorname{Re}}}
\renewcommand{\Im}{{\operatorname{Im}}}


%categories
\newcommand{\Grp}{{\textbf{Grp}}}
\newcommand{\Fld}{{\textbf{Fld}}}
\newcommand{\Sh}{{\textbf{Sh}}}
\newcommand{\Set}{{\textbf{Set}}}
\newcommand{\Mod}{{\textbf{Mod}}}
\newcommand{\PSh}{{\textbf{PSh}}}
\newcommand{\Sch}{{\textbf{Sch}}}
\newcommand{\Top}{{\textbf{Top}}}
\newcommand{\RgSp}{{\textbf{RgSp}}}
\newcommand{\Ring}{{\textbf{Ring}}}


%operatorname -- 3+ letters (lowercase and upper case)
\newcommand{\rad}{{\operatorname{rad}}}
\newcommand{\sep}{{\operatorname{sep}}}
\newcommand{\res}{{\operatorname{res}}}
\newcommand{\sgn}{{\operatorname{sgn}}}
\newcommand{\pre}{{\operatorname{pre}}}
\newcommand{\ord}{{\operatorname{ord}}}
\newcommand{\vol}{{\operatorname{vol}}}
\newcommand{\lcm}{{\operatorname{lcm}}}
\renewcommand{\div}{{\operatorname{div}}}

\newcommand{\conj}{{\operatorname{conj}}}
\newcommand{\disc}{{\operatorname{disc}}}
\newcommand{\stab}{{\operatorname{stab}}}
\newcommand{\kdim}{{\operatorname{kdim}}}
\newcommand{\diag}{{\operatorname{diag}}}

\newcommand{\trdeg}{\operatorname{trdeg}}
\newcommand{\coker}{{\operatorname{coker}}}
\newcommand{\codim}{{\operatorname{codim}}}

\newcommand{\Ann}{{\operatorname{Ann}}}
\newcommand{\Aut}{{\operatorname{Aut}}}
\newcommand{\Der}{{\operatorname{Der}}}
\newcommand{\Div}{{\operatorname{Div}}}
\newcommand{\Emb}{{\operatorname{Emb}}}
\newcommand{\End}{{\operatorname{End}}}
\newcommand{\Ext}{{\operatorname{Ext}}}
\newcommand{\Gal}{{\operatorname{Gal}}}
\newcommand{\Hom}{{\operatorname{Hom}}}
\newcommand{\Ind}{{\operatorname{Ind}}}
\newcommand{\Inn}{{\operatorname{Inn}}}
\newcommand{\Int}{{\operatorname{int}}}
\newcommand{\Irr}{{\operatorname{Irr}}}
\newcommand{\Jac}{{\operatorname{Jac}}}
\newcommand{\Mor}[1]{\operatorname{Mor}(\textbf{#1})}
\newcommand{\Nil}{{\operatorname{Nil}}}
\newcommand{\Obj}{{\operatorname{Obj}}}
\newcommand{\Out}{{\operatorname{Out}}}
\newcommand{\Pic}{{\operatorname{Pic}\,}}
\newcommand{\Rep}{{\operatorname{Rep}}}
\newcommand{\Res}{{\operatorname{Res}}}
\newcommand{\Spl}{{\operatorname{Spl}}}
\newcommand{\Syl}{{\operatorname{Syl}}}
\newcommand{\Sym}{{\operatorname{Sym}}}
\newcommand{\Tor}{{\operatorname{Tor}}}

\newcommand{\Char}{{\operatorname{char}}}
\newcommand{\PDiv}{{\operatorname{PDiv}}}
\newcommand{\Frac}{{\operatorname{Frac}}}
\newcommand{\Frob}{{\operatorname{Frob}}}
\newcommand{\Proj}{{\operatorname{Proj}\,}}
\newcommand{\RMod}{{}_R\operatorname{Mod}}
\newcommand{\SExt}{{\mathcal{E}\emph{xt}}}
\newcommand{\SHom}{{\emph{Hom}}}
\newcommand{\Span}{{\operatorname{span}}}
\newcommand{\Spec}{{\operatorname{Spec}\,}}
\newcommand{\Supp}{{\operatorname{Supp}}}
\newcommand{\Tens}[1]{#1\otimes --}
\newcommand{\fHom}[1]{\operatorname{Hom}(#1, --)}

\newcommand{\mSpec}{{\operatorname{mSpec}}}


% connectors
% \renewcommand{\mid}{\big|}
\newcommand{\sse}{\subseteq}
\newcommand{\subg}{\leqslant}
\newcommand{\supg}{\geqslant}
\newcommand{\ssne}{\subsetneq}
\newcommand{\isosubg}{\lesssim}
\newcommand{\nsse}{\not\subseteq}
\newcommand{\nsubg}{\trianglelefteq}
\newcommand{\nsupg}{\trianglerighteq}
\newcommand{\csubg}{\blacktriangleleft}
\renewcommand{\mid}{\big|}


% Arrows
\newcommand{\rw}{\rightarrow}
\newcommand{\Rw}{\Rightarrow}
\newcommand{\lw}{\leftarrow}
\newcommand{\Lw}{\Leftarrow}
\newcommand{\lrw}{\leftrightarrow}
\newcommand{\LRw}{\Leftrightarrow}
\newcommand{\acton}{\curvearrowright}
\newcommand{\actson}{\curvearrowright}
\newcommand\mapsfrom{\mathrel{\reflectbox{\ensuremath{\mapsto}}}}

% arrows for commutative diagram: create four new angled double-struck arrows
\newcommand\myrot[1]{\mathrel{\rotatebox[origin=c]{#1}{$\Rightarrow$}}}
\newcommand\NEarrow{\myrot{45}}
\newcommand\NWarrow{\myrot{135}}
\newcommand\SWarrow{\myrot{-135}}
\newcommand\SEarrow{\myrot{-45}}


%shorthands
\newcommand{\oln}[1]{{\overline{#1}}}
\renewcommand{\bf}[1]{\textbf{#1}}
\newcommand{\qn}{\quad\newline}


% limits
\newcommand{\Lim}{{\varprojlim}}
\newcommand{\coLim}{{\varinjlim}}
\newcommand{\Dlim}{\lim\limits_{\longrightarrow}}
\newcommand{\coDlim}{\lim\limits_{\longleftarrow}}
\newcommand{\cofHom}[1]{\operatorname{Hom}(--, #1)}


%for saturated ideals in the localization section (I don't like these, I think I'll delete)
\newcommand{\LIm}{{}^e }
\newcommand{\LPre}{{}^c }

%  ╔══════════════════════════════════════════════════════════╗
%  ║                    Custom Commands                       ║
%  ╚══════════════════════════════════════════════════════════╝

% swap varphi and phi
\let\temp\phi
\let\phi\varphi
\let\varphi\temp

\newcommand{\placeholder}{\_\_\_\_\_}

\newcommand{\set}[2]{\left\{#1 \ \middle|\ #2\right\}}

%mod should not have the crazy amount of space to its right!
\renewcommand{\mod}{\,\operatorname{mod}\,}

% dcups
\newcommand{\dcup}{\sqcup}
\newcommand{\Dcup}{\coprod}
% \newcommand{\Dcup}{\bigsqcup}

\newcommand{\nbot}{\mathrel{\rlap{$\bot$}\mkern2mu/}}

%a nice large circle
\newcommand{\tikzcircle}[2][red,fill=black]{\tikz[baseline=-0.5ex]\draw[#1,radius=#2] (0,0) circle ;}%

\makeatletter
\newcommand*\bigcdot{\mathpalette\bigcdot@{.5}}
\newcommand*\bigcdot@[2]{\mathbin{\vcenter{\hbox{\scalebox{#2}{$\m@th#1\bullet$}}}}}
\makeatother

%somehow not a default command
\def\rddots{\mathstrut^{.^{.^{.}}}}

%% new delimiter
\DeclarePairedDelimiter{\ceil}{\lceil}{\rceil}


%  ╔══════════════════════════════════════════════════════════╗
%  ║                     environments                         ║
%  ╚══════════════════════════════════════════════════════════╝


%remember the option: breakable

% create tcolor boxes
\newtcbtheorem[number within=section]{thm}{Theorem}%
{
colback=green!5,
colframe=green!35!black,
fonttitle=\bfseries,
label type=theorem
}{th}

\newtcbtheorem[number within=section, use counter from=thm]{lem}{Lemma}%
{
colback=blue!5,
colframe=black!35!black,
fonttitle=\bfseries,
label type=lemma
}{lm}
\newtcbtheorem[number within=section, use counter from=thm]{prop}{Proposition}%
{
colback=white!5,
colframe=white!35!black,
fonttitle=\bfseries,
label type=proposition
}{pr}
\newtcbtheorem[number within=section, use counter from=thm]{cor}{Corollary}%
{colback=blue!5,
colframe=blue!35!black,
fonttitle=\bfseries,
label type=corollary
}{co}
\newtcbtheorem[number within=section, use counter from=thm]{axiom}{Axiom}%
{colback=black!5,
colframe=yellow!35!black,
fonttitle=\bfseries,
label type=axiom
}{ax}
\newtcbtheorem[number within=section, use counter from=thm]{defn}{Definition}%
{colback=red!5,
colframe=red!35!black,
fonttitle=\bfseries,
label type=definition
}{df}


% for <s-cr> to create \item[$\LRw$] instead of \item
\newenvironment{equivEnumerate}
 {\begin{enumerate}
	}{
\end{enumerate}}

%insure the exercises are properly numbered
\counterwithin{Exercise}{section}
\counterwithin{Answer}{section}


%Custom enviroments
 \newenvironment{note}
 {\begin{tcolorbox}[enhanced, sharp corners, colback=white , borderline={0.2pt}{0pt}{black}]
   \begin{quote}\textbf{Note}:
   }{
   \end{quote}\end{tcolorbox}}

 \newenvironment{intuition}
 {\begin{quote}\textbf{Intuition}:
   }{
   \end{quote}}

\newtcbtheorem[number within=section, use counter from=thm]{titledBox}{}{%
  enhanced,
  sharp corners,
  colback=white,
  colframe=black,
  borderline={0.2pt}{0pt}{black},
  left=2mm,
  right=2mm,
  top=2mm,
  bottom=2mm,
  title style={opacity=1},
  attach title to upper,
  before upper={\textbf{\tcbtitletext}\quad},
  coltitle=black,
  fonttitle=\bfseries,
  separator sign={\quad}
}{box}


%better example and proof environments
\def\exampletext{Example} % If English
\colorlet{colexam}{red!55!black} % Global example color


\newtcbtheorem[number within=section, use counter from=thm]{example}{Example}{% \newtcolorbox[use counter=testexample]{testexamplebox}{%
% Example Frame Start
empty,% Empty previously set parameters
title={\exampletext \thetcbcounter #1},% use \thetcbcounter to access the testexample counter text
% Attaching a box requires an overlay
attach boxed title to top left,
   % Ensures proper line breaking in longer titles
   minipage boxed title,
% (boxed title style requires an overlay)
boxed title style={empty,size=minimal,toprule=0pt,top=4pt,left=3mm,overlay={}},
coltitle=colexam,fonttitle=\bfseries,
before=\par\medskip\noindent,parbox=false,boxsep=0pt,left=3mm,right=0mm,top=2pt,breakable,pad at break=0mm,
   before upper=\csname @totalleftmargin\endcsname0pt, % Use instead of parbox=true. This ensures parskip is inherited by box.
% Handles box when it exists on one page only
overlay unbroken={\draw[colexam,line width=.5pt] ([xshift=-0pt]title.north west) -- ([xshift=-0pt]frame.south west); },
% Handles multipage box: first page
overlay first={\draw[colexam,line width=.5pt] ([xshift=-0pt]title.north west) -- ([xshift=-0pt]frame.south west); },
% Handles multipage box: middle page
overlay middle={\draw[colexam,line width=.5pt] ([xshift=-0pt]frame.north west) -- ([xshift=-0pt]frame.south west); },
% Handles multipage box: last page
overlay last={\draw[colexam,line width=.5pt] ([xshift=-0pt]frame.north west) -- ([xshift=-0pt]frame.south west); }%
}{ex}



\def\prooftext{Proof} % If English
\NewDocumentEnvironment{Proof}{ O{} }
{
    \colorlet{colexam}{black!55!black} % Global example color
    \newtcolorbox[number within=section]{testproofbox}{%
        empty,% Empty previously set parameters
        title={\emph{\prooftext} \thetcbcounter: #1},
        attach boxed title to top left,
        minipage boxed title,
        boxed title style={empty,size=minimal,toprule=0pt,top=4pt,left=3mm,overlay={}},
        coltitle=colexam,
        fonttitle=\bfseries,
        before=\par\medskip\noindent,
        parbox=false,
        boxsep=0pt,
        left=3mm,
        right=0mm,
        top=2pt,
        breakable,
        pad at break=0mm,
        before upper=\csname @totalleftmargin\endcsname0pt,
        % Removed height constraints
        % Added flexible width
        width=\dimexpr\textwidth-3mm\relax,
        % Modified overlays
        overlay unbroken={\draw[colexam,line width=.5pt] ([xshift=-0pt]title.north west) -- ([xshift=-0pt]frame.south west); },
        overlay first={\draw[colexam,line width=.5pt] ([xshift=-0pt]title.north west) -- ([xshift=-0pt]frame.south west); },
        overlay middle={\draw[colexam,line width=.5pt] ([xshift=-0pt]frame.north west) -- ([xshift=-0pt]frame.south west); },
        overlay last={\draw[colexam,line width=.5pt] ([xshift=-0pt]frame.north west) -- ([xshift=-0pt]frame.south west); },%
    }
    \begin{testproofbox}
}
{\end{testproofbox}\endlist}


%  ╔══════════════════════════════════════════════════════════╗
%  ║                    for Dynkin diagrams                   ║
%  ╚══════════════════════════════════════════════════════════╝

\def\row#1/#2!{#1_{\IfStrEq{#2}{}{n}{#2}} & \dynkin{#1}{#2}\\}
\newcommand{\tble}[1]{
   \renewcommand*\do[1]{\row##1!}
   \[
      \begin{array}{ll}\docsvlist{#1}\end{array}
   \]
}

%  ╔══════════════════════════════════════════════════════════╗
%  ║                         links                            ║
%  ╚══════════════════════════════════════════════════════════╝

\definecolor{LinkColour}{HTML}{A64A3B} % favourite
% \definecolor{LinkColour}{HTML}{8C3D32} % 2nd place
\hypersetup{
  colorlinks,
  citecolor=black,
  filecolor=magenta,
  linkcolor=LinkColour,
  urlcolor=blue,
}



%  ╔══════════════════════════════════════════════════════════╗
%  ║                          misc                            ║
%  ╚══════════════════════════════════════════════════════════╝

\stackMath %to be able to stack text in math mode
\pgfplotsset{compat=1.18}
\makeindex

%import options: all, bibliography, book-formatting, booksSetting, customEnvironments, dynkin, hw-formatting, language-setup, newcommands, packages, preamble, proofs, settings, tcolorbox, test.svg,
\begin{document}
\maketitle
\tableofcontents

\textcolor{red}{I will try to have the introduction have as few prerequisites as possible so that what I'm
doing can at least be explainable, but it will be inevitable to bring up scheme theory and k-theory in the
paper so I will just drop the pretext of accessibility at that point}.

A natural question to ask in complex and algebraic geometry is ``how many
holomorphic/meromorphic/algebraic/analytic/etc functions are there on a (suitable) space $X$''. By letting
$\CF(X)$ represent our choice of functions (the global sections of a [quasi-coherent] sheaf $\CF$), we shall
denote this question a question about finding information about the $0$th sheaf cohomological group
\[
  H^0(X, \CF)
\]
which by~\cite{NateAlgGeo} is equal to $\CF(X)$. This group shall sometimes be trivial, and the reason
for its trivialness can be described by higher order sheaf cohomology groups:
\[
  H^1(X, \CF), H^2(X, \CF), ...
\]
exactly how to extract this information is known as the \emph{Riemann Roch Problem}. To motivate the direction
the problem took, let us answer this question in some simple cases.


The simplest case would be to deal with polynomial functions over $X = \C$ which are determined by a finite
set of roots (counting multiplicity) and a constant:
\[
  p(z) = c(z-a_1)\cdots (z-a_n)
\]
and similarly for rational functions. For future connection to sheaves let $\CO_X(X)$ represent all algebraic functions on $X$
and $\CM_X(X)$ represent all rational functions on $X$; other common notation is $\C[X] = \CO_X(X)$ and $\C(X)
= \CM_X(X)$ where $\C$ is often replaced with some other field $K$ or ring $R$. If we ask for holomorphic
functions so that $(X,\CF) = (\C, \CH(\C))$, then, then the Weierstrass Factorization Theorem
(see~\cite[section~4.2]{NateComplexAnalysis}) tells us that $f \in \CH(\C)$ factors as an infinite polynomial along with
some correction factors
\[
  z^m e^{g(z)}\prod_k^\infty \left[ \left( 1- \frac{z}{a_k} \right)e^{p_k(z)} \right]
\]
where
\[
  p_k(z) =  \frac{z}{a_k} + \frac{1}{2} \left( \frac{z}{a_k} \right)^2 + \cdots + \frac{1}{m_k} \left( \frac{z}{a_{m_k}} \right)^m
\]
is the correction factor and the $a_k$ are accumulation points; in particular there are analytic concerns.
A meromorphic function $f$ on $\C$ can be determined by two entire function $f = g/h$ (again
see~\cite[section~4.2]{NateComplexAnalysis}); this is possible as for poles of $f$ we can define a function $h$ with the
appropriate roots and multiplicities. In other words, determining meromorphic functions becomes
a classification of non-accumulating sequences in $\C$. The situation requires even more analytic information in hyperbolic
space $D^2$ (think of the existence of $\sum_n z^{n!}$ and $1/(1-z)$). To understand the algebraic behaviour, limiting the scope to
\emph{compact spaces} such as Riemann surfaces or projective varieties gets much better as any holomorphic/meromorphic
functions must have finitely many roots/poles by the Identity theorem (that if two functions agree on a set
of accumulation points, they are equal), and so we can focus our efforts within the algebraic setting. This
setting is plenty rich enough for study and often offers a good starting point to generalize to the
non-compact case\footnote{Though these cases are still very much being studied, see
\textcolor{red}{ref:HERE}}. In fact, by the efforts of Serre, Chow, and Grothendeick, the category of analytic
functions on (suitable) compact spaces is equivalent to the category of algebraic functions on these spaces
(\cite{hallGAGATheorems2022}, \cite{chowCompactComplexAnalytic1949}),
meaning compact spaces are ``intrinsically'' algebraic. We shall prove some of the pertinent parts of this
correspondence in \textcolor{red}{ref:HERE}.

Let us try finding holomorphic/meromorphic functions on $X = \CP^1$ and see if they are
determined by the points of $\CP^1$. Immediately, $\CP^1$ has only constant holomorphic functions by an easy
generalization of Liouville or the maximum modulus principle on manifolds. By the Residue Theorem, the sum of
the number of roots and poles of a meromorphic function must be $0$.

Let us represent this formally. Let
\[
  D = \sum_i n_ip_i \qquad p_i \in \CP^1
\]
be a finite formal sum; we may think of it as the free (abelian) group $F^{\Z}(\CP^1)$. Call $D$
a \emph{divisor}, and let let $\Div(\CP^1)$
represent this group. A divisor can be thought of as an object that remembers which roots/poles with which
multiplicity we are interested in. Define:
\[
  \div(f) =  \sum_i \nu_{p_i}(f)p_i
\]
where $\nu_{p_i}(f)$ is the order of the zero/pole, and is negative if it is a pole. This forms a subgroup of
$\Div(\CP^1)$; represent it as $\PDiv(\CP^1)$ (``P'' for principal), and elements of this subgroup are called
\emph{principal divisors}. This subgroup represents the divisors that are \emph{representable} as a
meromorphic function. Let us find which divisors can be represented. Let:
\[
  \deg D = \sum_i n_i
\]
Notice that $\deg(-)$ is certainly a group homomorphism. By our earlier remark:
\[
  \deg(\div(f)) = 0
\]
which is a restriction applied to any meromorphic function on $\CP^1$. It is easy to see that this restriciotn
would work on any compact Riemann surface (see~\cite[chapter~6]{NateComplexAnalysis}). If $\deg(D) = 0$, by our earlier remark
we have that there exists an $f$ such that $\div(f) = D$, showing this conditions characterizes meromorphic
functions $f \in \CM(\CP^1)$ (but it shall not on other Riemann surfaces). Now, consider the \emph{(Weil) divisor class group}
\[
  \Cl(\CP^1) = \frac{\Div(\CP^1)}{\PDiv(\CP^1)}
\]
for any $[D] \in \Cl(\CP^1)$,  $\deg([D])$ is well-defined as $\deg(D) = \deg(D + f)$ for any $f \in
\PDiv(\CP^1)$. If $\deg(D) = \deg(D')$, then $\deg(D) - \deg(D) = \deg(D-D') = 0$, and hence $D-D'
= \div(f)$. Therefore, we have that:
\[
  \Cl(\CP^1) \cong \Z
\]
which can be thought of as measuring the obstruction in a meromorphic function being defined by its
points/roots. As sheaf cohomology measures obstruction of local information having global information, we
shall have:
\[
  \Cl(\CP^1) \cong H^1(\CP^1, \CO_{\CP^1}^*)
\]
the details of this isomorphism and why $\CO_{\CP^1}^*$ shall be given in \textcolor{red}{ref:HERE}.


Translating the result for $(\C, \CM_\C)$, we see that:
\[
  \Cl(\C) = \frac{\Div(\C)}{\PDiv(\C)} \cong 0
\]
as every divisor can be represented by a principal divisor. We can interpret this as there being no
impediments to using local information to construct meromorphic functions, and we shall later also see that
all cohomology groups shall be trivial. Notice that the ring $\C[x]$ is a UFD; this algebraic fact shall be
given geometric meaning in \textcolor{red}{ref:HERE}.

Let us look at the next simplest nontrivial example. Let $X = \C/\Lambda$ where $\Lambda\sse \C$ is a lattice spanned by the
$\R$-linearly independent elements $(1,\tau)$ and let $\CF = \CM_X$ (all rational functions on $X$); $X$ is
called an \emph{elliptic curve}, and so we shall label it as $E$. As mentioned earlier, if $f \in \CM_E$,
$\div(f) = 0$. However, if $\deg(D) = 0$, this does not imply there exists an $f \in \CM_E$ such that $\div(f)
= D$. To see this, let analyze $\CM_E$ a little more closely. As $\C/\Lambda$ is the quotient by $\Lambda$,
any meromorphic function must respect the $\Lambda$ structure:
\[
  f(z+n+m\tau) = f(z) \qquad \forall n+m\tau \in \Lambda
\]
such a function is called \emph{doubly periodic} or \emph{elliptic function} (as they are defined on elliptic
curves). Using this periodic condition, it is simple to see that
$\div(f) = 0$ by looking at:
\[
  \frac{1}{2\pi i}\int_{\partial D} \frac{f'(z)}{f(z)}dz = 0
\]
where $D$ is the fundamental domain\footnote{If there are zero/poles on the path, we simply shift around
appropriately}, namely the opposing edges cancel out. We can modify this integral to find other properties,
for example:
\[
  \frac{1}{2\pi i}\int_{\partial D} z\frac{f'(z)}{f(z)}dz = n + m\tau \in \Lambda
\]
which is nonzero as $z$ is not doubly periodic. Replacing $z$ by other algebraic functions all end up in
$\Lambda$, and so we see that if $D = \div(f)$:
\[
  \sum_i n_ip_i \equiv 0 \mod \Lambda
\]
where here $\sum_i n_ip_i\notin F^\Z(E)$ but $\sum_i n_ip_i \in E$, namely we are trying $p_i$ as a complex
number multiplied by $n_i$ and summed. Certainly there is no such restriction on
a general divisor $D \in \Div(E)$ of degree $0$. We thus have shown that there is a surjective map:
\begin{equation}\label{eq:ellipticCrvIso}
  \Cl^0(E) = \frac{\Div^0(E)}{\PDiv(E)} \twoheadrightarrow \C/\Lambda
\end{equation}
where $\Div^0(E)$ are all divisor of degree $0$. It is in fact an isomorphism, which comes from constructing
an elliptic functions using the Weierstrass $\wp$-function, the details are
in~\cite[chapter~5]{NateComplexAnalysis}. Thus, we get the short exact sequence:
\[
  0 \to \C/\Lambda \to \Div(E) \xrightarrow{\deg} \Z \to 0
\]
often called the \emph{Jacobi-Abel} exact sequence as $\C/\Lambda$ is a Jacobian variety, and a short exact sequence:
\[
  0 \to \Cl^0(E) \to \Cl(E) \xrightarrow{\deg} \Z \to 0
\]
that together fit into the commutative diagram:
\[
\begin{tikzcd}
0 \arrow[r] & \mathrm{Div}^0(E) \arrow[r, hook] \arrow[d, two heads, "\pi"] & \mathrm{Div}(E) \arrow[r, "\deg"] \arrow[d, two heads, "\pi"] & \mathbb{Z} \arrow[r] \arrow[d, equals] & 0 \\
0 \arrow[r] & \Cl^0(E) \arrow[r, hook] & \Cl(E) \arrow[r, "\deg"] & \mathbb{Z} \arrow[r] & 0
\end{tikzcd}
\]
Continuing on to other Riemann surfaces, we shall loose the isomorphism given in
equation~(\ref{eq:ellipticCrvIso}) and is the beginning of the study of Jacobian
varieties (see~\cite[section~6.4]{NateComplexAnalysis}), while going in
higher dimensions is the beginning of the study of Weil and Cartier Divisors (see~\cite{NateAlgGeo}), and
study how far away a space is from being constructed by ``principal'' objects. Overall, the question becomes
less about finding which functions exist and more about some global properties of the space $X$.

To answer the Riemann Roch problem, a more granular approach to find which functions (of various kinds) exist
on a space $X$ must be taken. As divisors keep track of root/pole information, they are a great tool in
classifying functions that ``satisfy'' the properties of the divisor. To make this precise, define an order on
divisors $\div(X)$ where (for now) $X$ is either a Riemann surface or a smooth projective curve. Define
a partial order on $\Div(X)$ by saying $D_1
\ge D_2$ if each coefficient of $D_2 - D_1$ is non-negative. Then we can consider:
\[
  \CO_X(D)(X) = \set{f \in \CM_X(X)^*}{\div(f) + D \ge 0}
\]
In other words, if $D =\sum_i n_ip_i$, then $\CO_X(D)(X)$ is all meromorphic functions $f$ where it has a pole at $p_i$ of
order at most $n_i$ if $n_i > 0$ and roots at $p_i$ of order at least $|n_j|$ if $n_j < 0$ (and $\CO_X(D)$ is
a [coherent] sheaf). From this, we can
reformulate the Riemann Roch problem to a question about $H^0(X, \CO_X(D))$. The famous Riemann-Roch theorem as
proven by Riemann and Roch (where Riemann proved the Riemann-inequality, and Roch provided the
necessary work to make it an equality, see~\cite{a.griffithsIntroductionAlgebraicCurves}) is now known as
\emph{Riemann Roch for curves} and is the equality:
\begin{equation}\label{eq:riemannRoch1dIntro}
  \dim H^0(X, \CO_X(D)) - \dim H^1(X, \CO_X(D)) = \deg(D)+1-g
\end{equation}
where $g$ is the genus of the Riemann surface or smooth projective curve
(see~\cite{hartshorneAlgebraicGeometry2013} or~\cite{NateAlgGeo} for defining the genus of a curve via scheme
theory, or for a
direct computational proof using normalization see~\cite{NateComplexAnalysis}). In the original phrasing of the Riemann Roch
theorem, the group $H^1(X, \CO_X(D))$ is replaced with  $H^0(X, \CO(K_X-D))^*$ where $K_X$ is the canonical
divisor of $X$; these groups are isomorphic by Serre Duality (see~\cite[section~6.3]{NateAlgGeo}). Then this
space is identified with the space of meromorphic $1$-forms (either through some simple techniques from
complex analysis, see Brill-Noether Reciprocity in either~\cite{a.griffithsIntroductionAlgebraicCurves}
or~\cite{NateComplexAnalysis}, or by combining of Poincar\'e duality with DeRham's Theorem, see
either~\cite{leeIntroductionRiemannianManifolds2018} or~\cite{NateDiffGeo}). Switching to the traditional
notation $l(D) = \dim H^0(X, \CO_X(D))$ and $i(D) = \dim H^0(X, \CO_X(K_X-D))^*$ we get:
\[
  l(D) - l(K-D) = \deg(D) + 1-g
\]
\textcolor{red}{idk if I will prove RR for curves or relegate it to a book}. As we phrased the Riemann-Roch
problem in the beginning, this formula would then be asking for the value of $l(D)$. There are cases where
$l(K-D)$ disappears (for example if $D = 0$ or $D = K$), and there are many cases we can compute. It is simple
to compute the right hand side but in general it is harder to computer the left hand side. Thus, this formula
becomes a new constriction that helps us find $l(D)$. There are a set of theorems known as \emph{vanishing
theorems} that specify under which conditions $l(K-D)$ (or more generally, higher cohomology groups) vanish
(see~\cite{deligneRelevementsModulop2etDecomposition1987}); under these conditions we exact recover $l(D)$.


Let us work towards the intuition behind the generalization. The notation on left hand side in
equation~(\ref{eq:riemannRoch1dIntro}) can be compactified by using the Euler-Poincar\'e characteristic (or
simply the Euler characteristic): $\chi(X, \CO_X(D)) = \sum_{i=0}^\infty (-1)^i \dim H^i(X, \CO_X(D))$ so
that:
\[
  \chi(X, \CO_X(D)) = \deg(D) + 1-g
\]
Note that $H^i(X, \CO_X(D)) = 0$ for $i > 1$, \textcolor{red}{cite? Reference for later?} giving the equality.
The right hand side shall be drastically different, split into special invariants of vector bundles known as
the Chern class and Todd class \textcolor{red}{mention the ``duality'' between them? see notes} . In fact, the
Euler Characteristic shall also turn out to be an invariant on vector bundles, and hence the Riemann Roch
theorem generalizes to a question about the existence of sections on vector bundles over a compact manifolds.
We shall mention the generalization to smooth quasi-projective schemes over a field near the end, and also
talk a bit about a recent paper (\cite{navarroRiemannRochFormulaProjective2017}) generalizing to finite
dimensional Noetherian schemes.

To motivate how vector bundles will enter the picture, it is worth-while seeing how $\CO_X(D)$ can be
interpreted as a line bundle. To that end, recall that the projective curve $X = \P^1$ as defined in classical
algebraic geometry only has constant global sections. Namely, if $U_0, U_1\cong \A^1$ is the usual cover with
global sections (i.e. functions) $\CO_X(U_0) \cong \CO_X(U_1) = K[x]$ and the compatibility condition on
$U_0\cap U_1$ is $f(z) = g(1/z)$, namely:
\begin{center}
  \[
\begin{tikzcd}
U_0                                               &               U_1 \\
     U_0\cap U_1 \arrow[r, "\cong", no head] \arrow[u, hook] & U_0\cap U_1 \arrow[u, hook]
\end{tikzcd}
\qquad \text{i.e.} \qquad
\begin{tikzcd}
\A^1                                       &              \A^1 \\
     \A^1\setminus\{0\} \arrow[u, hook] \arrow[r, "\cong", no head] & \A^1\setminus\{0\} \arrow[u, hook]
\end{tikzcd}
\]
\end{center}
gives:
\begin{center}
  \[
\begin{tikzcd}
\CO_X(U_0) \arrow[d]                                &              \CO_X(U_1) \arrow[d] \\
                    \CO_X(U_0\cap U_1) \arrow[r, "\cong", no head] & \CO_X(U_0\cap U_1)
\end{tikzcd}
\qquad \text{i.e.} \qquad
\begin{tikzcd}
{K[x]} \arrow[d]                        & {K[y]} \arrow[d] \\
{K[x, 1/x]} \arrow[r, "\cong", no head] & {K[y, 1/y]}
\end{tikzcd}
\]
\end{center}
where $x \mapsto 1/y$ and $1/x \mapsto y$, Then if $F([x':y']) \in
\CO_{\P^1}(\P^1)$ is a global function, then $F([x:1]) \to f(x) \in \CO_X(U_0)$ and $F([1:y]) \to g(y) \in \CO_X(U_1)$
are the restrictions of the function to $U_0$ and $U_1$ respectively, and by the above diagram it must be that
these functions map to the same element in the restriction, namely:
\[
  f(z) \to f(x) = f(1/y) = g(y) \leftarrow g(y)
\]
but it is impossible that $f(1/y) = g(y)$ unless $f,g$ are the same constant. Thus, it must be that the only
global functions on the projective line are constants:
\[
  \CO_{\P^1}(\P^1) = K
\]
The fundamental limitation here is that the restriction $f(x) = g(1/x)$ is very strong. What if we can loosen
this restriction? For example what if
\[
  f(x) = xg\left(\frac{1}{x}\right)
\]
Then we can in fact form linear homogeneous functions! For example consider
\[
  F([x':y']) = ax'+by'
\]
Then $F$ maps to $f(x) = ax+b$ and $g(y) = a + by$. Then:
\[
  f(x) = x\left(a + b \frac{1}{x}\right) = ax+b = xg\left(\frac{1}{x}\right)
\]
which works! More generally, if the condition $f(x) = x^ng(1/x)$ allows for degree $n$ homogeneous
polynomials. Denote the degree $n$-homogeneous polynomials by $\CO_{\P^1}(n)(\P^1)$.



\textcolor{red}{also mention how the weil/cartier divisor stuff earlier isomorphic to line bundles, hence that
discussion really about properties of line bundles}.

% . Notice that they can be determined by taking some countable set of non-accumulating
% of points of $\C$. For any proper simply connected open set $U \sse \C$ we have that it is
% biholomorphic to the open unit disk $D^2$ by the Riemann mapping theorem, and so all holomorphic functions
% with power series representations at $0$ that converge with radius $< 1$ are in $\CH(D)$; this includes
% $\frac{1}{1-z}$ and pathological functions like $\sum_n z^{n!}$. In other words, these are much more
% difficult and requires some more tools from analysis.


\newpage
\printbibliography
\end{document}
